% source: David Mumford, "Abelian Varieties", TIFR Studies in Mathematics 5, Oxford UP, 1970
% pdf_page: 0061
% doc_page: 50 (Chapter II, Section 5, "Cohomology and base change")
% retrieved: 2026-07-11
% transcribed_by: claude-fable-5 (vision, rendered at 200dpi; scanned book, no text layer)

\DeclareMathOperator{\Image}{Im}

% [running head: ABELIAN VARIETIES, page 50]

% [this Corollary is printed unnumbered; later proofs (pp. 52, 54) refer to it as Corollary 1]
\paragraph{Corollary.}
\textit{Let $X$, $Y$, $f$ and $\mathcal{F}$ be as in the theorem (except
that $Y$ need not be affine). Then we have:}
\begin{itemize}
  \item[(a)] \textit{For each $p \geq 0$, the function $Y \to \mathbf{Z}$
    defined by \\
    $y \longmapsto \dim_{k(y)} H^p(X_y, \mathcal{F}_y)$ is upper
    semicontinuous on $Y$.}
  \item[(b)] \textit{The function $Y \to \mathbf{Z}$ defined by}
    \[
      y \to \chi(\mathcal{F}_y)
        = \sum_{p=0}^{\infty} (-1)^p \dim_{k(y)} H^p(X_y, \mathcal{F}_y)
    \]
    \textit{is locally constant on $Y$.}
\end{itemize}

\paragraph{Proof.}
The problem being local on $Y$, we may assume $Y$ affine. Let $K^{\cdot}$ be
a complex as in the proposition; by further localization, we may assume
$K^{\cdot}$ to be a free complex. Denote by $d^p : K^p \to K^{p+1}$ the
coboundary of $K$. We then have
\[
\begin{aligned}
  \dim_{k(y)} H^p(X_y, \mathcal{F}_y)
    &= \dim_{k(y)}[\ker(d^p \otimes_A k(y))] - {} \\
    &\qquad\qquad - \dim_{k(y)}[\Image(d^{p-1} \otimes_A k(y))] \\
    &= \dim_{k(y)}[K^p \otimes k(y)]
       - \dim_{k(y)}[\Image(d^p \otimes k(y))] - {} \\
    &\qquad\qquad - \dim_{k(y)}[\Image(d^{p-1} \otimes k(y))].
       \qquad (*)
\end{aligned}
\]
The first term being constant on $Y$, (b) follows on taking alternating sum
of $(*)$ over all $p$. We assert that for any $p \geq 0$, the function
$\rho_p(y) = \dim_{k(y)}[\Image(d^p \otimes k(y))]$ is lower
semi-continuous on $Y$. In fact, if $r$ is any integer $\geq 0$, and
$d^p_r : \Lambda^r K^p \to \Lambda^r K^{p+1}$ is the map induced by $d^p$,
\[
  \{\, y \in Y \mid \rho_p(y) < r \,\}
    = \{\, y \in Y \mid d^p_r \otimes k(y) = 0 \,\},
\]
and this set is closed since $d^p_r$ is a homomorphism of free finitely
generated modules, and hence is described by a matrix in $A$, and the above
set is the set of common zeros of all entries of the matrix. This proves
(a).

Moreover, the theorem gives the following criterion for putting together
the cohomology groups of $\mathcal{F}$ along the fibres of $f$ into a
vector bundle on $Y$.

\paragraph{Corollary 2.}
\textit{Let $X$, $Y$, $f$ and $\mathcal{F}$ be as above. Assume $Y$ is
reduced and connected. Then for all $p$ the following are equivalent:}
% [statement of Corollary 2 continues on the next page]
